%================================================================%
%=====================  Modelo de Monografia ====================%
%========================  UFOP - DEPRO =========================%
% Proposta de texto em conformidade com normas da ABNT ----------%
% implementadas pelo projeto abntex2, que pode ser acessado pela %
% página  http://abntex2.googlecode.com/  -----------------------%
%================================================================%
%========================== AUTORES =============================%
% Marcelus Xavier Oliveira  email: marcelusxavier@gmail.com =====%
% Dayanne Gouveia Coelho  email: dayagc@gmail.com ===============%
%================================================================%
%======================== IMPORTANTE ============================%
% Para o bem comum dos estudantes e professores do DEPRO, é      %
% importante mantermos apenas um template de monografia. Assim,  %
% para qualquer sugestão no template, favor entrar em contato com%
% os autores ----------------------------------------------------% 
%================================================================%
%======================== Versão 2.0 ============================%
%================================================================%
\documentclass[12pt, % tamanho da fonte
	%openright,	% capítulos começam em pág ímpar
	oneside,		  % twoside para impressão em frente e verso.  
	a4paper,			% tamanho do papel. 
	english,			% Idioma adicional para hifenização
	brazil				% Idioma principal 
	]{abntex2}
	
\include{structure}  % Estrutura do Texto e Pacotes Principais

% -- Informações para Capa e Folha de Rosto: ---------------
\titulo{Título da Monografia}
\autor{Nome Aluno} \autorcite{Aluno,Nome}
\local{Ouro Preto} \uf{MG}
\data{03 de novembro de 2017} \ano{2017}
\orientador{Prof. Dr. Nome do Orientador}  % Nome do orientador 
\ttorientador{Universidade Federal de Ouro Preto} % Instituição do orientador
\coorientador{Prof. Dr. Nome do Coorientador}   % Nome do coorientador
\ttcoorientador{Universidade Federal de Ouro Preto} % Instituição do Coorientador
\instituicao{Universidade Federal de Ouro Preto} \sigla{UFOP}
\instituto{Escola de Minas}
\departamento{Departamento de Engenharia de Produção}
\curso{Engenharia de Produção}	
\tipotrabalho{Monografia}
\grau{Bacharel em Engenharia de Produção}

%------Nomes dos examinadores.  
\examinadorum{Prof. Dr. Membro da Banca 1} \ttexaminadorum{Universidade Federal de ... - UFXX}
\examinadordois{Prof. Dr. Membro da Banca  2} \ttexaminadordois{Universidade Federal de ... - UFXX}
%\examinadortres{Prof. Dr. Membro da Banca  3} \ttexaminadortres{Universidade Federal de ... - UFXX}
%\examinadorquatro{Prof. Dr. Membro da Banca  4} \ttexaminadorquatro{Universidade Federal de ... - UFXX}

% ------------------------------------------------------
\makeindex   

\begin{document} % Início do documento

\frenchspacing  % Retira espaço obsoleto entre as frases.

% ----------------------------------------------------------
% -- Elementos Pré-Textuais: -------------------------------
\pagenumbering{roman}

\imprimircapa  % Capa
\imprimirfolhaderosto % Folha de rosto

\include{PreTextuais/FichaCatalografica}
\include{PreTextuais/FichaAprovacao} 
\include{PreTextuais/Dedicatoria}
\include{PreTextuais/Agradecimento}
\include{PreTextuais/Epigrafe}
\include{PreTextuais/Resumo}

\renewcommand{\listfigurename}{Lista de figuras}
\pdfbookmark[0]{\listfigurename}{lof}
\listoffigures*   % Cria a Lista de Figuras
\cleardoublepage

\pdfbookmark[0]{\listtablename}{lot}
\listoftables*  % Cria a lista de Tabelas
\cleardoublepage

\renewcommand{\listalgorithmcfname}{Lista de algoritmos}
\pdfbookmark[0]{\listalgorithmcfname}{lof}
\listofalgorithmes   % Cria a lista de Tabelas
\cleardoublepage

\include{PreTextuais/ListaSiglas}
\include{PreTextuais/ListaSimbolos}


\pdfbookmark[0]{\contentsname}{toc}
\tableofcontents*
\cleardoublepage




% ----------------------------------------------------------
% -- Capítulos do Trabalho: --------------------------------
\pagenumbering{arabic} \setcounter{page}{1}
\textual 
\include{Capitulos/Cap1-Introducao}
\include{Capitulos/Cap2-RevisaoBibliografica}
\include{Capitulos/Cap5-Resultados}
\include{Capitulos/Cap6-Conclusao}



% ----------------------------------------------------------
% -- Elementos Pós-Textuais: -------------------------------
\postextual  
\bibliography{bibliografia} % Referências bibliográficas
\include{PosTextuais/Apendice}
\include{PosTextuais/Anexo}
%\phantompart  \printindex  % Indice Remissivo
% ----------------------------------------------------------
\end{document}  % fim do documento
